%% sections/05-discussion.tex

\section{Discussion}
\label{sec:discussion}

\subsection{Replicability of the Mechanism}

The innovation-cluster-driven amendment mechanism requires three infrastructure components:
(a) a session log format that captures specific outcomes with sufficient textual detail
to support evidence citation; (b) a systematic logging protocol that distinguishes
innovations from ordinary session events; and (c) a forward-only versioning convention
for the rubric being amended.

None of these components is specific to TRPG evaluation. Any creative evaluation
pipeline that produces session-level artifacts, uses multi-dimensional scoring, and
encounters quality outcomes that resist existing anchors can implement this mechanism.
The 2-innovation threshold is conservative: it prevents single-instance patterns from
driving amendments while allowing rapid response to genuinely recurring phenomena.
In the Marathon corpus, the lowest cluster size that produced an adopted amendment
was 2 (13 of the 15 clusters had size 2 or 3; only the AL v1.9 cluster had 4
constituent innovations).

\subsection{The Atmospheric Landing Asymmetry}

Atmospheric Landing (AL) received 8 of the 13 amendments, far more than any other
dimension. This asymmetry has a structural explanation: AL is the dimension that most
directly scores complex, emergent, inter-PC narrative phenomena. Mechanical Fairness
(3 amendments) scores properties of the dice system; Module Fidelity (2 amendments)
scores coverage of designed content; Character Agency (2 amendments) scores
choice-making behavior. All three have relatively bounded quality spaces. AL, by
contrast, encompasses the full range of ways that emotional and atmospheric intent
can be realized in session play --- a space that is effectively unbounded.

The implication for rubric design is that high-dimensional evaluations of complex
narrative systems should expect the ``emotional realization'' dimension to require
the most ongoing amendment work. Designers of such systems should budget for this
ongoing work rather than treating it as a sign of rubric failure.

\subsection{Archetype-Driven Amendment Bursts}

The C4 reactivation --- 4 amendments from a single campaign after 2 campaigns produced
only 2 each --- suggests that campaign archetype is a stronger predictor of amendment
bursts than session count alone. The siege-resource mechanic introduced genuine quality
phenomena not encountered in grief-adjacent campaigns: a PC managing spell slots
across 7 sessions as a unified resource (grief-paragraph resource expression), the
intra-session restraint of placing a token beside tools rather than showing the party,
the directed foreknowledge of a Portent die voiced to a named ally before the roll.

These are not subtle variations on existing patterns; they are structurally different
quality achievements enabled by the resource-exhaustion mechanic. The rubric had no
anchor for them because the previous campaigns had not produced a sustained-pressure
context in which they could occur.

This finding has a methodological implication: a rubric evaluated on a corpus with
only one or two campaign archetypes will be under-specified for others. The Marathon
corpus's 7 distinct archetypes is a minimum adequate sample for producing a
reasonably generalized rubric; we would expect 2--3 additional amendments from each
new archetype introduced into the pipeline.

\subsection{Limitations}

The mechanism requires a reviewer who is both familiar with the rubric and alert to
outcomes that exceed it. In the Marathon corpus, the AI system performing gate review
also performed innovation logging. This raises the question of whether the system's
own cognitive patterns (its implicit theories of what constitutes quality) shaped
what was logged as an innovation versus what was absorbed into an existing anchor
without logging. We cannot rule out systematic biases in the innovation log as a result.

Additionally, the corpus is entirely from a single research environment. The innovation
rate in production TRPG sessions --- with human players, real-time play, and more
varied adventure designs --- might differ substantially. We do not claim that 120+
innovations over 49 sessions is a universal rate; we claim only that the mechanism
will produce a declining-and-reactivating innovation rate as the rubric matures
and novel archetypes are introduced.
